% 杨舒云自用LaTex笔记设定集
%下面是一些碎碎念……


% ---------------------------------------------------------------------
% 自定义颜色（流萤主题配色——老婆的颜色！）
\usepackage{xcolor}
\definecolor{fblack}{HTML}{3E324A} % #3e324a（紫黑）
\definecolor{fgrayblue}{HTML}{475D7B} % #475d7b（灰蓝）
\definecolor{fgraygreen}{HTML}{97C6C0} % #97c6c0（灰绿）
\definecolor{fred}{HTML}{E26E1B} % #e26e1b（深橘红）
\definecolor{fsilver}{HTML}{E6E4E0} % #e6e4e0（银白）
\definecolor{fbluegreen}{HTML}{4DF8E8} % #4df8e8（蓝绿）


% ---------------------------------------------------------------------
% 链接
\usepackage{ctex}
\usepackage[top=28mm,bottom=28mm,left=15mm,right=15mm]{geometry}
\usepackage{hyperref} 
\hypersetup{
	colorlinks,
	linktoc = section, % 超链接位置，选项有section, page, all
	linkcolor = fgrayblue, % linkcolor 目录颜色
	citecolor = fgrayblue  % citecolor 引用颜色
}


% ---------------------------------------------------------------------
% 引入必要的包
%A
\usepackage{amsmath,enumerate,multirow,float} % 数学公式，
\usepackage{tabularx}
\usepackage{tabu}
\usepackage{subfig}
\usepackage{fancyhdr}
\usepackage{graphicx} % 插图
\usepackage{wrapfig}  
\usepackage{physics}
\usepackage{appendix}
\usepackage{amsfonts}
%B
\usepackage{mathrsfs} % 字体
\usepackage{calligra}
\usepackage{lipsum}
\usepackage{adjustbox}% 调整表格大小
\usepackage{tabularray} % 绘制表格时可以更加方便添加框线
%C
\usepackage{caption}
\usepackage{authblk}
\usepackage{braket}
\usepackage{amssymb} % 包含更多数学符号
\usepackage[en-US]{datetime2} % 设置日期为美式英语格式
\usepackage{geometry}  % 页面调整


% ---------------------------------------------------------------------
% 对目录、章节标题的设置
\usepackage{titlesec}
\usepackage{titletoc}
% 目录
\renewcommand{\contentsname}{\centerline{\Huge Outline}}

% 定义章节标题格式
\titleformat{\chapter}[block]
{\centering\bfseries\color{fblack}\Huge} % 格式：居中、加粗、Large字体
{PART \thechapter}          % 标题前缀
{1em}                       % 前缀与标题之间的间距
{ }                         % 标题格式
% 设置章节编号格式为罗马数字
\renewcommand{\thechapter}{\Roman{chapter}}

% 自定义章节编号
\renewcommand\thesection{\S\arabic{section}}
\renewcommand\thesubsection{\thesection.\arabic{subsection}}
\renewcommand\thesubsubsection{\S\arabic{section}.\arabic{subsection}.\arabic{subsubsection}}


% 自定义section格式
\titleformat{\section}
{\normalfont\bfseries\color{fgrayblue}\Large} % 设置字体、大小、加粗、颜色
{\thesection} % 设置section编号的格式
{1em} % 编号和标题之间的间距
{}

% 自定义subsection格式
\titleformat{\subsection}
{\normalfont\bfseries\color{fgrayblue}\large} % 设置字体、大小、颜色
{\thesubsection} % 设置subsection编号的格式
{1em} % 编号和标题之间的间距
{}

% 自定义subsubsection格式
\setcounter{secnumdepth}{3}
\titleformat{\subsubsection}
{\normalfont\bfseries\color{fgrayblue}\normalsize} % 设置字体、大小、颜色
{\thesubsubsection} % 设置subsubsection编号的格式
{1em} % 编号和标题之间的间距
{}

% 章节计数器独立于Chapter
\makeatletter
\@addtoreset{section}{chapter}
\@addtoreset{subsection}{chapter}
\@addtoreset{subsubsection}{chapter}
\makeatother


% ---------------------------------------------------------------------
% 颜色盒子
\usepackage{tcolorbox}
\tcbuselibrary{skins, breakable}
% 设置 Key Point 环境
\newtcolorbox{keypoint}[1][]{enhanced, breakable,
	colback=fbluegreen!30!white,colframe=fbluegreen!80!black,colbacktitle = fbluegreen!80!black,
	attach boxed title to top left = {yshift = -2mm, xshift = 5mm},
	boxed title style = {sharp corners},
	fonttitle=\bfseries, title=\textbf{Key Point} \thechapter.\arabic{keypointcounter}, #1}

\newcounter{keypointcounter}[chapter]
\renewcommand{\thekeypointcounter}{\thechapter.\arabic{keypointcounter}}

% 修改引用的格式
\AtBeginEnvironment{keypoint}{\refstepcounter{keypointcounter}}

% 设置highlight环境
\newtcbox{\highlight}[1][fbluegreen]
{on line, arc = 0pt, outer arc = 0pt,
	colback = #1, colframe = white,
	boxsep = 0pt, left = 1pt, right = 1pt, top = 2pt, bottom = 2pt,
	boxrule = 0pt, bottomrule = 1pt, toprule = 1pt}
	
% 设置一般框架ubox（universal box）环境，以及一些不错的调整方案
\newtcolorbox{ubox}[2][]
{enhanced, breakable,
	colback = white, colframe = fbluegreen, coltitle = black,colbacktitle = fbluegreen,
	attach boxed title to top left = {yshift = -2mm, xshift = 5mm},
	boxed title style = {sharp corners},
	fonttitle = \bfseries,
	title={#2},#1}
	
% 设置macbox框架
\definecolor{wt1}{HTML}{ebebeb} % 白色主题的标题框顶部背景颜色
\definecolor{wt2}{HTML}{bebebe} % 白色主题的标题框底部背景颜色
\definecolor{wt3}{HTML}{efefef} % 白色主题的正文部分的背景颜色
\definecolor{circ1}{HTML}{eb605b} % 标题框左边第一个圆圈的颜色
\definecolor{circ2}{HTML}{f6bb31} % 标题框左边第二个圆圈的颜色
\definecolor{circ3}{HTML}{56cb45} % 标题框左边第三个圆圈的颜色
\newtcolorbox{wmacbox}[2][]
{enhanced, breakable,
	boxrule=0pt,
	frame code={
		\clip[rounded corners=1mm] (frame.south west) rectangle (frame.north east);
		\shade[top color=wt1,bottom color=wt2] (title.south west) rectangle (title.north east);
		\fill[circ1] ([xshift=  2em]title.west) circle [radius=1ex];
		\fill[circ2] ([xshift=3.5em]title.west) circle [radius=1ex];
		\fill[circ3] ([xshift=  5em]title.west) circle [radius=1ex];
	},
	halign title=flush center,
	colback = wt3,
	fonttitle = \bfseries,
	title=\textcolor{fblack}{{#2}},
	#1}


% ---------------------------------------------------------------------
% 利用cleveref改变引用格式，\cref是引用命令
\usepackage{cleveref}
% 设置引用格式
\crefformat{figure}{#2\textbf{\textcolor{fred}{figure~#1}}#3}
\crefformat{table}{#2\textbf{\textcolor{fred}{table~#1}}#3}
\crefformat{equation}{#2\textbf{\textcolor{fred}{formula~(#1)}}#3}
\crefname{keypointcounter}{Key Point}{\textcolor{fred}{key point}}
\Crefname{keypointcounter}{Key Point}{\textcolor{fred}{key point}}
\crefformat{keypointcounter}{#2\textbf{\textcolor{fred}{key-point~#1}}#3}

% ---------------------------------------------------------------------
%	页眉页脚设置
\fancypagestyle{plain}{\pagestyle{fancy}}
\pagestyle{fancy}
\lhead{Physics with Coding} % 左边页眉
%\chead{\nouppercase{\leftmark}} % 这个还没搞好，但是我不想搞了
\rhead{\kaishu Copy@pifuyuini} % 右边页眉
\cfoot{\thepage} % 页脚，中间添加页码


% ---------------------------------------------------------------------
% 列表的设置
\usepackage{enumitem}
% 设置有序列表格式
%\setlist[enumerate,1]{label=\textcolor{fbluegreen}{\arabic*.}, font=\bfseries\color{fbluegreen}}
%\setlist[enumerate,2]{label=\textcolor{fbluegreen}{(\arabic*)}, font=\bfseries\color{fbluegreen}}
%\setlist[enumerate,3]{label=\textcolor{fbluegreen}{\roman*.}, font=\bfseries\color{fbluegreen}}
%\setlist[enumerate,4]{label=\textcolor{fbluegreen}{(\roman*)}, font=\bfseries\color{fbluegreen}}

% 自定义列表格式
\setlist[enumerate,1]{label=\textbf{\textcolor{fbluegreen}{\arabic*.}}, font=\large\bfseries\color{fbluegreen}}
\setlist[enumerate,2]{label=\textbf{\textcolor{fbluegreen}{(\arabic*)}}, font=\large\bfseries\color{fbluegreen}}
\setlist[enumerate,3]{label=\textbf{\textcolor{fbluegreen}{\roman*.}}, font=\large\bfseries\color{fbluegreen}}
\setlist[enumerate,4]{label=\textbf{\textcolor{fbluegreen}{(\roman*)}}, font=\large\bfseries\color{fbluegreen}}


% 设置无序列表格式
\setlist[itemize,1]{label=\textcolor{fbluegreen}{$\blacktriangleright$}, font=\bfseries\color{fbluegreen}}
\setlist[itemize,2]{label=\textcolor{fbluegreen}{$\bullet$}, font=\bfseries\color{fbluegreen}}
\setlist[itemize,3]{label=\textcolor{fbluegreen}{$\blacksquare$}, font=\bfseries\color{fbluegreen}}


% ---------------------------------------------------------------------
% 图表的设置
% 设置图表和表格编号格式
\renewcommand{\thefigure}{\thechapter.\arabic{figure}}
\renewcommand{\thetable}{\thechapter.\arabic{table}}

% 设置图表和表格标题格式
\captionsetup[figure]{
	labelfont={color=fgrayblue,bf}, % 设置Figure X.Y的颜色
	labelsep=period,
	name=Figure
}

\captionsetup[table]{
	labelfont={color=fgrayblue,bf}, % 设置Table X.Y的颜色
	labelsep=period,
	name=Table
}


% ---------------------------------------------------------------------
% 公式的设置（报错就不用）
\usepackage{etoolbox} % 用于修改公式标签颜色
% 设置公式编号格式
\newcounter{eqncounter}[section]
\renewcommand{\theeqncounter}{\thechapter.\thesection.\arabic{eqncounter}}

% 使用 etoolbox 修改公式标签颜色
\makeatletter
\preto\equation{\colorlet{saved@eqcolor}{.}\color{fgrayblue}}
\appto\equation{\color{saved@eqcolor}}
\def\tagform@#1{\maketag@@@{\color{fgrayblue}(#1)}}
\makeatother

% 自动编号公式
\newcommand{\numeq}[1]{\refstepcounter{eqncounter}\label{#1}\tag{\theeqncounter}}


% ---------------------------------------------------------------------
% 制作一个代码块代，分为两个部分：外层框架用tcolorbox，内层文本展示用listing

% 相比于前面的wmacbox，这个macbox可以适配代码展示listing，但我仍在前面保留了那个不适配的作为参考，二者的制作方式略有差异，但颜色我就直接套用前面的了。
%\definecolor{macosbox@top}{RGB}{237,237,237}
%\definecolor{macosbox@bot}{RGB}{189,189,189}
\definecolor{macosbox@bord}{RGB}{182,176,176}
%\definecolor{macosbox@bg}{RGB}{240,240,240}
\newtcolorbox{macbox}[2][]{
	enhanced,
	breakable,
	coltitle=black,
	colback = wt3,%macosbox@bg,
	boxrule=0mm,
	frame style={draw=macosbox@bord,fill=macosbox@bord},
	title style={top color=wt1,bottom color=wt2},
	drop fuzzy shadow=black,
	title={{\textcolor{circ1}{\huge$\bullet$}
			\textcolor{circ2}{\huge$\bullet$}
			\textcolor{circ3}{\huge$\bullet$}
			\hspace*{\fill}\texttt{#2}\hspace*{77mm}\hspace*{\fill}}},#1
}

%   listing代码环境设置
\usepackage{listings}
\lstloadlanguages{python}
\lstdefinestyle{pythonstyle}{
	%backgroundcolor=\color{gray!6},
	language=python,
	frameround=tftt,
	%frame=shadowbox, 
	keepspaces=true,
	breaklines,
	columns=spaceflexible,      
	basicstyle=\ttfamily\color{fblack}, % 基本文本设置（这里可以改改字号）
	keywordstyle=[1]\color{fbluegreen!90!black}\bfseries, 
	keywordstyle=[2]\color{fred!90!black},   
	stringstyle=\color{Purple!80!fblack},       
	showstringspaces=false,
	commentstyle=\ttfamily\color{fgraygreen!65!black},% 注释文本设置
	tabsize=2,
	morekeywords={as},
	morekeywords=[2]{np, plt, sp},
	numbers=left, % 代码行数
	numberstyle=\scriptsize\color{fgrayblue}, % 代码行数的数字字体设置
	stepnumber=1,
	numbersep=0pt % 代码行号和代码主体之间的距离
	%rulesepcolor=\color{fsilver}
}


% ---------------------------------------------------------------------
%	其他设置
\def\degree{${}^{\circ}$} % 角度
\graphicspath{{./images/}} % 插入图片的相对路径
\allowdisplaybreaks[4]  %允许公式跨页